\documentclass[10pt,aspectratio=169]{beamer}

\usetheme{Madrid}
\usecolortheme{default}
\setbeamertemplate{navigation symbols}{}
\setbeamertemplate{footline}[frame number]

\usepackage{amsmath,amssymb,amsthm}
\usepackage{mathtools}
\usepackage{bm}
\usepackage{booktabs}
\usepackage{colortbl}
\usepackage{graphicx}

\newcommand{\R}{\mathbb{R}}
\newcommand{\N}{\mathbb{N}}
\newcommand{\calX}{\mathcal{X}}
\newcommand{\bA}{\mathbf{A}}
\newcommand{\bP}{\mathbf{P}}
\newcommand{\bI}{\mathbf{I}}
\newcommand{\bM}{\mathbf{M}}
\newcommand{\bp}{\mathbf{p}}
\newcommand{\br}{\mathbf{r}}
\newcommand{\bH}{\mathbf{H}}
\newcommand{\btheta}{\bm{\theta}}
\newcommand{\bphi}{\bm{\phi}}
\newcommand{\bzeta}{\bm{\zeta}}
\newcommand{\cL}{\mathcal{L}}

\title[CME Inference + Koopman]{Sparse CME Inference with Fr\'{e}chet Optimization\\and Koopman Diagnostics}
\subtitle{Initialization, Optimization, and Spectral Quality Checks}
\author{}
\date{}

\begin{document}

\begin{frame}
\titlepage
\end{frame}

\begin{frame}{Talk Map}
\textbf{Central question:} when should we trust inferred CME generators?
\begin{itemize}
    \item \textbf{Part I (inference core):} initialization and Fr\'{e}chet-STLS optimization.
    \item \textbf{Part II (evidence):} quantitative accuracy and distribution-level fidelity.
    \item \textbf{Part III (Koopman audit):} spectral conditioning, pseudoinverse sensitivity, and stability.
\end{itemize}

\vspace{0.4em}
\begin{itemize}
    \item connect CME-rate recovery to operator-theoretic diagnostics that can motivate a SIADS narrative.
\end{itemize}
\end{frame}

\begin{frame}{Inverse Problem}
\textbf{CME on finite state space} $\calX=\{x_1,\ldots,x_n\}$:
\[
\dot{\bp}(t)=\bA(\btheta)\bp(t),\qquad \bp(t+\Delta t)=e^{\bA(\btheta)\Delta t}\bp(t)
\]

\vspace{0.3em}
\textbf{Mass-action parameterization:}
\[
a_m(x)=\btheta_m^\top \bphi(x),\qquad
\bA(\btheta)=\sum_{m,k}\theta_{m,k}\,\bA_m^{(k)}
\]

\vspace{0.3em}
\textbf{Objective:}
\[
\min_{\btheta\ge 0}\; \cL(\btheta)=\frac12\sum_{t=0}^{T-1}
\left\|\hat{\bp}_{t+1}-e^{\bA(\btheta)\Delta t}\hat{\bp}_{t}\right\|_2^2
\]
with sparsity enforced by STLS thresholding.
\end{frame}

\begin{frame}{Initialization Regimes}
\textbf{When diagonal dominance is plausible}:
\[
\Delta t \le \frac{\log 2}{w_{\max}}\;\Rightarrow\;
\hat{\bA}=\frac{1}{\Delta t}\log\!\left(\hat{\bP}^{\Delta t}\right)
\]

\vspace{0.4em}
\textbf{Robust alternative (always valid CME generator)}:
\[
\hat{\bA}_{\text{lin}}=\frac{\hat{\bP}^{\Delta t}-\bI}{\Delta t}
\]
Properties:
\begin{itemize}
    \item off-diagonals nonnegative
    \item column sums exactly zero
\end{itemize}

\vspace{0.3em}
\textbf{Bias (clean paper):}
\[
\hat{\bA}_{\text{lin}}=\bA^*+\frac{\Delta t}{2}(\bA^*)^2+O(\Delta t^2)
\]
Refinement stage removes most of this initialization bias.
\end{frame}

\begin{frame}{Fr\'{e}chet Gradient Engine}
\textbf{Fr\'{e}chet derivative of matrix exponential:}
\[
L_{\exp}(\bM,\bH)=\int_0^1 e^{(1-s)\bM}\bH e^{s\bM}\,ds
\]

\vspace{0.3em}
\textbf{Block identity:}
\[
\exp\!\begin{pmatrix}\bM&\bH\\0&\bM\end{pmatrix}
=
\begin{pmatrix}e^{\bM}&L_{\exp}(\bM,\bH)\\0&e^{\bM}\end{pmatrix}
\]

\vspace{0.3em}
\textbf{Objective gradient:}
\[
\frac{\partial \cL}{\partial \theta_{m,k}}
=
-\Delta t\sum_t
\br_t^\top
L_{\exp}\!\left(\bA\Delta t,\bA_m^{(k)}\right)\hat{\bp}_t
\]

\vspace{0.3em}
Krylov implementation gives per-gradient complexity:
\[
O\!\left(T\cdot M\cdot K\cdot m_K\cdot \mathrm{nnz}(\bA)\right),\quad m_K\approx 20\!-\!30.
\]
\end{frame}

\begin{frame}{Optimization Loop (What We Keep)}
\textbf{Only the relevant loop: initialization + Fr\'{e}chet-STLS}
\begin{enumerate}
    \item Build initial rates from sampled transitions (empirical exposure-weighted estimate).
    \item Run L-BFGS updates on $\cL(\btheta)$ using exact Fr\'{e}chet gradients.
    \item Project to nonnegative rates.
    \item Hard-threshold small coefficients for sparsity.
    \item Polish on discovered support.
\end{enumerate}

\vspace{0.4em}
\textbf{Theory used here}
\begin{itemize}
    \item local support recovery (under restricted Jacobian regularity)
    \item statistical rate: $\|\hat{\btheta}-\btheta^*\|_2=O\!\left(\sqrt{s^*\log(MK)/T}\right)$
\end{itemize}
\end{frame}

\begin{frame}{ Discussion}
\begin{itemize}
    \item Initialization gives a physically valid first generator even when matrix-log assumptions fail.
    \item Fr\'{e}chet gradients make the nonlinear matrix-exponential fit numerically tractable at scale.
    \item STLS separates \textit{structure discovery} from \textit{coefficient refinement}.
\end{itemize}

\vspace{0.4em}
\textbf{Expected behavior from theory}
\begin{itemize}
    \item high-probability support recovery when the restricted Jacobian is well-conditioned;
    \item error decay with sample size on the true support.
\end{itemize}
\end{frame}

\begin{frame}{Most Recent Benchmark Summary}
\begin{center}
\begin{tabular}{lccc}
\toprule
& \textbf{Lotka-Volterra} & \textbf{Michaelis-Menten} & \textbf{Brusselator} \\
\midrule
Trajectories $N$ & 1000 & 1000 & 400 \\
Sampling interval & 0.1 & 0.1 & 0.2 \\
State-space size & 1652 & 58 & 309 \\
\midrule
Empirical init. error & 13.3\% & 69.5\% & 60.6\% \\
\rowcolor[gray]{0.9} Final error & \textbf{1.6\%} & \textbf{3.1\%} & \textbf{3.0\%} \\
\rowcolor[gray]{0.9} Error reduction & 8.5$\times$ & 22$\times$ & 20$\times$ \\
Structure recovery & 100\% & 100\% & 100\% \\
\bottomrule
\end{tabular}
\end{center}

\vspace{0.3em}
\textbf{Interpretation:}
\begin{itemize}
    \item Initialization can be strongly biased (especially MM and Brusselator), but Fr\'{e}chet-STLS consistently closes the gap.
    \item Exact structure recovery across all three systems indicates the sparse mechanism is identified correctly.
\end{itemize}
\end{frame}

\begin{frame}{MM Distribution Check}
\begin{center}
\includegraphics[width=0.95\textwidth]{figures/fig_distributions.pdf}
\end{center}
\vspace{-0.2em}
\textbf{Interpretation:}
\begin{itemize}
    \item Predicted and true probability mass align over time, so the model captures full transient law evolution.
    \item This is stronger than parameter-fit agreement: it validates the inferred semigroup action on distributions.
\end{itemize}
\end{frame}

\begin{frame}{STLS Stage Behavior}
\begin{center}
\includegraphics[width=0.95\textwidth]{figures/fig_stls_stages.pdf}
\end{center}
\vspace{-0.2em}
\textbf{Interpretation:}
\begin{itemize}
    \item Stage-wise curves separate support identification from numerical polishing.
    \item Most quantitative gain happens after Fr\'{e}chet optimization on the selected support.
\end{itemize}
\end{frame}

\begin{frame}{Convergence and Statistical Scaling}
\begin{center}
\includegraphics[width=0.92\textwidth]{figures/convergence_combined.pdf}
\end{center}
\vspace{-0.2em}
\textbf{Interpretation:}
\begin{itemize}
    \item Error decreases with data and optimization progress, matching the expected finite-sample trend.
    \item No instability plateau is visible in the tested regimes.
\end{itemize}
\end{frame}

\begin{frame}{Realistic Pipeline Comparison}
\begin{center}
\includegraphics[width=0.92\textwidth]{figures/realistic_pipeline_comparison.pdf}
\end{center}
\vspace{-0.2em}
\textbf{Interpretation:}
\begin{itemize}
    \item The full pipeline remains accurate under practical sampling limits and model complexity.
    \item Improvement over initialization persists in realistic, non-idealized settings.
\end{itemize}
\end{frame}

\begin{frame}{Koopman Diagnostics: What We Added}
\textbf{Goal:} audit quality of empirical transition operators before/after filtering.

\vspace{0.3em}
\textbf{Preconditioning used:}
\begin{itemize}
    \item keep states with obs count $\ge 50$ and $\max_t p_i(t)\ge 10^{-4}$
    \item cap pruning at 60\%
    \item account for dropped-state leakage on diagonal
\end{itemize}

\vspace{0.3em}
\textbf{Diagnostics:}
\begin{itemize}
    \item spectra of empirical $\bP$ and recovered generator $\bA_{\text{rec}}$
    \item singular spectrum and pseudoinverse gains
    \item Penrose residuals + $\|\bP^+\|$ norms
\end{itemize}

\vspace{0.3em}
\textbf{Why this matters for SIADS framing:}
\begin{itemize}
    \item it ties statistical inference quality to operator conditioning and mode reliability.
\end{itemize}
\end{frame}

\begin{frame}{Koopman Quantitative Summary (Preconditioned)}
\begin{center}
\begin{tabular}{lcc}
\toprule
\textbf{Metric} & \textbf{Lotka-Volterra} & \textbf{Brusselator} \\
\midrule
Kept states & 636/1587 & 98/245 \\
$\kappa(\bP)$ before $\to$ after & $2.84\!\times\!10^{12}\to 5.76\!\times\!10^7$ & $1.84\!\times\!10^{13}\to 4.92\!\times\!10^3$ \\
$\|\bP^+\|_2$ & $5.36\times 10^7$ & $4.03\times 10^3$ \\
Penrose residual scale & $10^{-9}$ to $10^{-11}$ & $10^{-13}$ to $10^{-15}$ \\
Unstable eigs of $\bA_{\text{rec}}$ & 0 & 0 \\
\bottomrule
\end{tabular}
\end{center}

\vspace{0.3em}
\textbf{Interpretation:}
\begin{itemize}
    \item conditioning improves strongly in both systems
    \item Brusselator remains well-conditioned post-inversion
    \item LV still has stronger inverse sensitivity in weakly sampled modes
\end{itemize}
\end{frame}

\begin{frame}{Lotka-Volterra: Transition Modes and Generator Spectrum}
\begin{center}
\includegraphics[width=0.80\textwidth,height=0.50\textheight,keepaspectratio]{figures/koopman_transition_modes_lotka_volterra.png}
\end{center}
\textbf{Interpretation:}
\begin{itemize}
    \item Transition eigenvalues concentrate near the real axis and close to $1$, indicating slow mixing directions.
    \item Recovered generator eigenvalues remain in the stable half-plane, but spread indicates sensitivity in weakly sampled regions.
\end{itemize}
\end{frame}

\begin{frame}{Lotka-Volterra: Pseudoinverse Diagnostics}
\begin{center}
\includegraphics[width=0.80\textwidth,height=0.50\textheight,keepaspectratio]{figures/koopman_pseudoinverse_lotka_volterra.png}
\end{center}
\textbf{Interpretation:}
\begin{itemize}
    \item Large row/column norm spikes in $\bP^+$ identify localized amplification channels.
    \item Penrose residuals are small, so the issue is not algebraic correctness but inverse sensitivity.
\end{itemize}
\end{frame}

\begin{frame}{Brusselator: Transition Modes and Generator Spectrum}
\begin{center}
\includegraphics[width=0.80\textwidth,height=0.50\textheight,keepaspectratio]{figures/koopman_transition_modes_brusselator.png}
\end{center}
\textbf{Interpretation:}
\begin{itemize}
    \item Modes are concentrated on a smaller active manifold with clearer spectral separation.
    \item Generator spectrum is stable and less dispersed, consistent with better conditioning.
\end{itemize}
\end{frame}

\begin{frame}{Brusselator: Pseudoinverse Diagnostics}
\begin{center}
\includegraphics[width=0.80\textwidth,height=0.50\textheight,keepaspectratio]{figures/koopman_pseudoinverse_brusselator.png}
\end{center}
\textbf{Interpretation:}
\begin{itemize}
    \item $\sigma(\bP)$ decays more gently after filtering and $\|\bP^+\|$ remains moderate.
    \item This explains why recovered generator modes are numerically more robust than in Lotka-Volterra.
\end{itemize}
\end{frame}

\begin{frame}{Cross-System Koopman Discussion}
\textbf{What the diagnostics say about inference reliability}
\begin{itemize}
    \item Both systems are spectrally stable after preconditioning (no unstable recovered generator eigenvalues).
    \item Conditioning gains are dramatic, but residual inverse sensitivity differs by system geometry and sampling support.
    \item Koopman-side diagnostics provide a practical confidence score for when recovered rates should be trusted.
\end{itemize}

\vspace{0.4em}
\textbf{Open technical question for discussion}
\begin{itemize}
    \item can we jointly optimize rate inference and state-space filtering to control both bias and operator conditioning?
\end{itemize}
\end{frame}

\begin{frame}{Meeting Summary and Discussion Prompts}
\textbf{Core message}
\begin{itemize}
    \item Fr\'{e}chet-based sparse optimization gives accurate CME recovery across three systems.
    \item Robust initialization is essential; Fr\'{e}chet-STLS provides the accuracy lift.
    \item Koopman diagnostics give an operator-level health check (conditioning, inverse sensitivity, spectral stability).
\end{itemize}

\vspace{0.4em}
\textbf{Technical caveat}
\begin{itemize}
    \item preconditioning improves stability but can bias rates if aggressive in high-leakage regions (seen more in LV than Brusselator).
\end{itemize}

\vspace{0.4em}
\textbf{Discussion prompts for Igor Mezic}
\begin{itemize}
    \item threshold/leakage sweep to choose a Pareto point: minimal bias with acceptable $\kappa(\bP)$ and $\|\bP^+\|$;
    \item operator-theoretic regularization ideas for ill-conditioned empirical transition matrices;
    \item positioning relative to existing GeDMD work: finite-state generator inference with conditioning-aware diagnostics.
\end{itemize}
\end{frame}

\end{document}
